% \iffalse meta-comment
%
% Copyright (C) 2017 Vincent Goulet, David Beauchemin
%
% This file may be distributed and/or modified under the conditions
% of the LaTeX Project Public License, either version 1.3c of this
% license or (at your option) any later version. The latest version
% of this license is in:
%
%   http://www.latex-project.org/lppl.txt
%
% and version 1.3c or later is part of all distributions of LaTeX
% version 2006/05/20 or later.
%
% This work has the LPPL maintenance status `maintained'.
%
% The Current Maintainer of this work is Vincent Goulet
% <vincent.goulet@act.ulaval.ca>.
%
% This work consists of the files actuarialsymbol.dtx and
% actuarialsymbol.ins and the derived files listed in the README file.
%
% \fi
%
% \iffalse
%<*dtx>
\ProvidesFile{actuarialsymbol.dtx}
%</dtx>
%<package>\NeedsTeXFormat{LaTeX2e}[2009/09/24]
%<package>\ProvidesPackage{actuarialsymbol}[2017/03/15 v0.1 Actuarial symbol]
%
%<*driver>
\documentclass[11pt,x11names,english]{ltxdoc}
  \usepackage[utf8]{inputenc}
  \usepackage[T1]{fontenc}
  \usepackage{natbib}
  \usepackage{babel}
  \usepackage{microtype}
  \usepackage[scaled=0.92]{helvet}
  \usepackage[scaled=1.02]{inconsolata}
  \usepackage[sc]{mathpazo}
  \usepackage{fontawesome}
  \usepackage{actuarialsymbol}
  \usepackage{graphicx,xcolor,framed}
  \usepackage{enumitem,tabularx,booktabs,amsthm}
  \DisableCrossrefs
  \CodelineNumbered
  \RecordChanges
  \GlossaryPrologue{\section*{Historique des versions}%
    \addcontentsline{toc}{section}{Historique des versions}}

  \definecolor{link}{rgb}{0,0.4,0.6}   % ~RoyalBlue de dvips
  \definecolor{url}{rgb}{0.6,0,0}      % rouge foncé
  \definecolor{citation}{rgb}{0,0.5,0} % vert foncé
  \colorlet{shadecolor}{LightYellow1} % vert foncé

  %% Liste description alignée à gauche
  \setlist[description]{leftmargin=*,align=left}

  \usepackage{hyperref}
  \hypersetup{colorlinks, linktocpage,
    urlcolor=url, linkcolor=link, citecolor=citation,
    bookmarksopen, bookmarksnumbered, bookmarksdepth=subsubsection,
    pdftitle={Manuel de référence de la classe ulthese},
    pdfauthor={Faculté des études supérieures et postdoctorales}}
  \addto\extrasfrench{%
    \def\appendixautorefname{annexe}%
    \def\tableautorefname{tableau}%
    \def\subsectionautorefname{section}%
  }

  \newenvironment{demo}{%
    \begin{trivlist}\item}{%
    \end{trivlist}}
  \newenvironment{texample}[1][0.5]{%
    \noindent\begin{minipage}{#1\linewidth}\begin{snugshade*}%
      \def\producing{\end{snugshade*}\end{minipage}\hfill\begin{minipage}{\dimexpr0.97\linewidth-#1\linewidth}%
        \hbox\bgroup\kern-.2pt%
        \vbox\bgroup\parindent0pt\relax
        % The 3pt is to cancel the -\lineskip from \displ@y
        \abovedisplayskip3pt \abovedisplayshortskip\abovedisplayskip
        \belowdisplayskip0pt \belowdisplayshortskip\belowdisplayskip
        \noindent}
    }{%
      \par
      % Ensure that a lonely \[\] structure doesn't take up width less than
      % \hsize.
      \hrule height0pt width\hsize
      \egroup\kern-.2pt\egroup
    \end{minipage}%
    \par
  }

  \theoremstyle{definition}
  \newtheorem*{rem}{Remark}

  \newcommand{\class}[1]{\textsf{#1}}
  \newcommand{\pkg}[1]{\textbf{#1}}

  \newcommand{\link}[2]{\href{#1}{#2~\raisebox{-0.2ex}{\faExternalLink}}}
  \newcommand{\doc}[3][documentation]{%
    \link{#3}{#1}\marginpar{\hfill\faBookmark~\texttt{#2}}}

\begin{document}
  \DocInput{actuarialsymbol.dtx}
\end{document}
%</driver>
% \fi
% \CheckSum{0}
%
% \GetFileInfo{actuarialsymbol.dtx}
% \title{Composition of actuarial symbols}
% \author{Vincent Goulet \\ École d'actuariat, Université
%   Laval \and David Beauchemin}
% \date{}
% \maketitle
%
% \begin{abstract}
%
% \end{abstract}
%
% \section{Introduction}
% \label{sec:introduction}
%
% Actuaries denote various quantities of life contingencies like
% present values of life insurances and life annuities, annual
% premiums, or reserves using a whole array of symbols. The resulting
% highly descriptive yet compact notation was standardized in as far
% back as 1898 \citep{Wolthuis:notation:2004}. \autoref{fig:mosaic}
% shows a creative use of the notation by the graduating class
% of 1972 in Actuarial Science at Université Laval.
%
% \begin{figure}[t]
%   \centering
%   \includegraphics[width=\linewidth]{mosaic}
%   \caption{``Actuariat'' (French for Actuarial Science) written
%   using actuarial symbols on the 1972 graduating class mosaic at
%   Université Laval}
%   \label{fig:mosaic}
% \end{figure}
%
% \citet[Appendix~4]{Bowers:2e:2002} offer an excellent overview of
% the composition rules for symbols of actuarial functions. In a
% nutshell, a core symbol, say $S$, is combined with auxiliary symbols
% positioned in subscript or in superscript, to the left or to the
% right. Schematically, we thus have:
% \begin{displaymath}
%   \actsymb[\fbox{\small I}\!\;][\fbox{\small II}\!\;]%
%      {\framebox[5ex]{\large $S$}}{\,\fbox{\small III}}%
%      [\,\fbox{\small IV}]
% \end{displaymath}
%
% The core symbol is in general a single letter. The letter may be
% ``accented'' with a bar ($\bar{A}$), double dots ($\ddot{a}$) or a
% circle ($\mathring{e}$). When the core symbol consists of two
% letters, they are grouped between parentheses, as in $\IA$ or
% $\Dadd$. Most commonly, there is an auxiliary symbol (or collection
% of symbols) in the lower-right position \fbox{III}. Otherwise,
% auxiliary symbols appear lower-left \fbox{I}, upper-left
% \fbox{II} and upper-right \fbox{IV}, in that order of frequency.
%
% Principle symbols for benefit premiums, reserves and amount of
% reduced paid-up insurance, $P$, $V$ and $W$ are combined with
% benefit symbols unless the benefit is a level unit insurance payable
% at the end of the year of death. In such cases, we have the
% following symbol structure (replace $P$ by $V$ or $W$ as needed):
% \begin{displaymath}
%   \actsymb[\fbox{\small I}\!\;][\fbox{\small II}\!\;]%
%      {\framebox[5ex]{\large $P$}}{}[\,\fbox{\small IV}]%
%   \bigl(\actsymb{\framebox[5ex]{\large $S$}}{\,\fbox{\small III}}\bigr)%
% \end{displaymath}
%
% Perhaps the most commonly used auxiliary symbol not readily
% available in {\LaTeX} is the ``angle'' denoting a duration $n$, as
% in $\angln$. Package \pkg{actuarialangle} \citep{actuarialangle}
% provides commands to create this symbol. This package is imported at
% load time by \pkg{actuarialsymbol}.
%
% Package \pkg{actuarialsymbol} provides a generic command to position
% all subscripts and superscripts easily and consistently around an
% actuarial symbol, and a few other commands to achieve some
% specialized typesetting. The package also defines a number of
% shortcuts to create the most common actuarial functions of financial
% mathematics and life contingencies.
%
% \section{Package features}
% \label{sec:features}
%
% We first describe the generic commands provided by
% \pkg{actuarialsymbol} to compose actuarial symbols.
%
% \begin{DescribeMacro}{\actsymb}
%   The generic command \cmd{\actsymb} typesets a core symbol with
%   surrounding subscripts and superscripts. Its syntax is somewhat
%   unusual for {\LaTeX}, but it serves well the natural order of the
%   building blocks of a symbol and their relative prevalence:
%   \begin{quote}
%     \cmd{\actsymb}\oarg{ll}\oarg{ul}\marg{core}\marg{lr}\oarg{ur}
%   \end{quote}
%   Above, \meta{ll} identifies the lower left subscript \fbox{I}
%   (following the notation in the schematic representation of the
%   Introduction); \meta{ul} is the upper left superscript \fbox{II};
%   \meta{core} is the core symbol $S$; \meta{lr} is the lower right
%   subscript \fbox{III}; \meta{ur} is the upper right superscript
%   $\fbox{IV}$.  The core symbol and the right subscript are
%   required, the other arguments are optional.
%   \begin{demo}
%     \begin{texample}
%       |\actsymb{A}{x}|
%       \producing
%       $\actsymb{A}{x}$
%     \end{texample}
%     \begin{texample}
%       |\actsymb[n]{A}{x}|
%       \producing
%       $\actsymb[n]{A}{x}$
%     \end{texample}
%     \begin{texample}
%       |\actsymb[n][2]{A}{x}|
%       \producing
%       $\actsymb[n][2]{A}{x}$
%     \end{texample}
%     \begin{texample}
%       |\actsymb[n][2]{A}{x}[(m)]|
%       \producing
%       $\actsymb[n][2]{A}{x}[(m)]$
%     \end{texample}
%   \end{demo}
%
% \end{DescribeMacro}
%
% \begin{rem}
%   {\TeX} adjusts the position of a subscript downward when a
%   superscript is present:
%   \begin{displaymath}
%     A_x \quad A_x^2.
%   \end{displaymath}
%   The package maintains this behaviour, but ensures that the left
%   and right subscripts, when both present, are at the same height:
%   \begin{displaymath}
%     \actsymb[t]{A}{x} \quad \actsymb[t]{A}{x}[2].
%   \end{displaymath}
%   (This is something ad~hoc constructions like |{}_{t}A_x^2| do not
%   provide.) If you would prefer a uniform subscript position
%   \emph{throughout your document}, load package \pkg{subdepth}
%   \citep{subdepth}.
% \end{rem}
%
% \begin{DescribeMacro}{\twoletsymb}
%   Writing two-letter core symbols like $\DA$ as |$(DA)$| results in
%   letters that are too distant from one another: $(DA)$. To unify
%   presentation, the package provides the command
%   \begin{quote}
%     \cmd{\twoletsymb}\marg{symbol}
%   \end{quote}
%   to create a two-letter \meta{symbol} surrounded by parentheses.
%   The symbol is typeset using \cmd{\mathit} to ensure proper spacing
%   between the letters. The package defines control structures for
%   all the main two-letter actuarial symbols; see
%   \autoref{sec:shortcuts}.
% \end{DescribeMacro}
%
% \begin{DescribeMacro}{\nthtop}
%   \begin{DescribeMacro}{\varnthtop}
%     Commands
%     \begin{quote}
%       \cmd{\nthtop}\oarg{mult}\marg{number}\marg{item} \\
%       \cmd{\varnthtop}\oarg{mult}\marg{number}\marg{item}
%     \end{quote}
%     put a precedence \meta{number} above an \meta{item}, smashed so
%     that the apparent height of the item is its normal height. This
%     is normally used in the right subscript \fbox{III} of a symbol.
%     With \cmd{\nthtop}, the spacing between the precedence number
%     and the item is constant. This can result in precedence numbers
%     placed at different heights if one item contains an horizontal
%     rule:
%     \begin{demo}
%       \begin{texample}[0.85]
%         |\actsymb{A}{\nthtop{1}{x}y:\nthtop{2}{\angln}}|
%         \producing
%         \rule{0pt}{12pt}$\actsymb{A}{\nthtop{1}{x}y:\nthtop{2}{\angln}}$
%       \end{texample}
%     \end{demo}
%     Conversely, \cmd{\varnthtop} always leaves enough space for an
%     intervening horizontal rules, such that all precedence numbers
%     lie at the same height:
%     \begin{demo}
%       \begin{texample}[0.85]
%         |\actsymb{A}{\varnthtop{1}{x}y:\varnthtop{2}{\angln}}|
%         \producing
%         \rule{0pt}{12pt}$\actsymb{A}{\varnthtop{1}{x}y:\varnthtop{2}{\angln}}$
%       \end{texample}
%     \end{demo}
%     The optional argument \meta{mult} of both commands is a
%     multiplier governing the spacing between the number and the
%     symbol. (Default values are $2$ for \cmd{\nthtop} and $7$ for
%     \cmd{\varnthtop}; see \autoref{sec:implementation} for details.)
%     The package also defines shortcuts for first, second and third
%     precedence; see \autoref{sec:shortcuts}.
%   \end{DescribeMacro}
% \end{DescribeMacro}
%
% \begin{DescribeMacro}{\nthbottom}
%   Similar to the above two commands,
%   \begin{quote}
%     \cmd{\nthbottom}\marg{number}\marg{item}
%   \end{quote}
%   puts a precedence \marg{number} below an \marg{item}. The package
%   also defines shortcuts for first and second precedence; see
%   \autoref{sec:shortcuts}.
% \end{DescribeMacro}
%
% \begin{rem}
%   The fact that top precedence numbers have zero height means they
%   will clash with a right superscript \fbox{IV}:
%   \begin{demo}
%     \begin{texample}[0.7]
%       |\actsymb{A}{\nthtop{1}{x}:\angln}[(m)]|
%       \producing
%       $\actsymb{A}{\nthtop{1}{x}:\angln}[(m)]$
%     \end{texample}
%   \end{demo}
%   In such rare circumstances, one needs to insert a \emph{strut} (an
%   invisible vertical rule) in the subscript to move it down as
%   needed:
%   \begin{demo}
%     \begin{texample}[0.7]
%       |\newcommand{\str}{\rule{0pt}{2.3ex}}| \\
%       |\actsymb{A}{\str\nthtop{1}{x}:\angln}[(m)]|
%       \producing
%       $\actsymb{A}{\rule{0pt}{2.3ex}\nthtop{1}{x}:\angln}[(m)]$
%     \end{texample}
%   \end{demo}
%   This also applies to bottom precedence numbers in inline formulas.
% \end{rem}
%
%    \begin{macrocode}
%\iffalse
%<*package>
%\fi
\RequirePackage{amsmath}
\RequirePackage{actuarialangle}

\DeclareRobustCommand{\actsymb}{\ACTS@actsymb}
\newcommand\ACTS@actsymb{\@ifnextchar[{\ACTS@@actsymb}{\ACTS@@actsymb[]}}
\newcommand\ACTS@@actsymb{}
\def\ACTS@@actsymb[#1]{\@ifnextchar[{\ACTS@@@actsymb[#1]}{\ACTS@@@actsymb[#1][]}}
\newcommand\ACTS@@@actsymb{}
\def\ACTS@@@actsymb[#1][#2]{\@ifnextchar[{\ACTS@@@@actsymb[#1][#2]}{\ACTS@@@@actsymb[#1][#2][]}}
\newcommand\ACTS@@@@actsymb{}
\def\ACTS@@@@actsymb[#1][#2][#3]#4#5{\@ifnextchar[{\@actsymbol[#1][#2][#3]{#4}{#5}}{\@actsymbol[#1][#2][#3]{#4}{#5}[]}}

\def\@actsymbol[#1][#2][#3]#4#5[#6]{
  \setbox\z@\hbox{#3}
  \@mathmeasure\@ne\scriptstyle{#1}
  \@mathmeasure\tw@\scriptstyle{#2}
  \@mathmeasure\thr@@\scriptstyle{#5}
  \@mathmeasure4\scriptstyle{#6}
  \mathop{}
  %% adjust depth of left and right *supercripts*
  \ifdim \dp\tw@>\dp4
    \setbox4\hbox to\wd4{\raisebox{\z@}[\ht4][\dp\tw@]{\unhbox4}}
  \else
    \setbox\tw@\hbox to\wd\tw@{\raisebox{\z@}[\ht\tw@][\dp4]{\unhbox\tw@}}
  \fi
  %% adjust width of *left* subscript and superscript
  \ifdim \wd\tw@>\wd\@ne
    \setbox\@ne\hbox to\wd\tw@{\hfil\unhbox\@ne}
  \else
    \setbox\tw@\hbox to\wd\@ne{\hfil\unhbox\tw@}
  \fi
  %% adjust height of left and right *subscripts*
  \ifdim \ht\@ne>\ht\thr@@
    \setbox\thr@@\vbox to \ht\@ne{\vfil\hbox to\wd\thr@@{\unhbox\thr@@}}
  \else
    \setbox\@ne\vbox to \ht\thr@@{\vfil\hbox to\wd\@ne{\unhbox\@ne}}
  \fi
  \ifnum\ifdim \wd\tw@>\z@ 1\else\ifdim \wd4>\z@ 1\else 0\fi\fi
    =1%  expression with superscripts either left or right
    \mathopen{\vphantom{#4}}^{\box\tw@}\sb{\box\@ne}%
    \ifdim\wd\z@=\z@ #4^{\box4}\sb{\box\thr@@}\else #3^{\box4}(#4^{}\sb{\box\thr@@})\fi
  \else% expression without superscripts
    \vphantom{#4}\sb{\box\@ne}%
    \ifdim\wd\z@=\z@ #4\sb{\box\thr@@}\else #3(#4\sb{\box\thr@@})\fi
  \fi
}

\DeclareRobustCommand{\twoletsymb}[1]{(\mathit{#1})}

\DeclareRobustCommand{\nthtop}[3][2]{\mathpalette{\preced@ t0{#1}{#2}}{#3}}
\DeclareRobustCommand{\varnthtop}[3][7]{\mathpalette{\preced@ t{#1}{0}{#2}}{#3}}
\DeclareRobustCommand{\nthbottom}[2]{\mathpalette{\preced@ t02{#1}}{#2}}

% In \preced@, #1 is `t' for top or `b' for bottom, #2 is a
% multiplier for the baseline skip between the top and bottom symbols, #3 is a
% multiplier for the line skip between the top and bottom symbols, #4
% is the first argument from the user, #5 is \displaystyle or
% \textstyle or \scriptstyle or \scriptscriptstyle, from \mathpalette,
% and #6 is the second argument given by the user. This peculiar
% ordering of the arguments is done to work around the restriction of
% \mathpalette that it only reads two arguments.

\def\preced@#1#2#3#4#5#6{%
  % Measure the arguments:
  \setbox\tw@\hbox{$\m@th#5#4$}%
  \setbox\z@\hbox{$\m@th#5#6$}
  \dimen@\wd\z@
  \vbox to\ht\z@{%
    % \lineskip is set using AMSTeX's \ex@, if available, for a slight
    % refinement in the spacing if this macro is used in eightpoint
    % text. If \ex@ is not available, \p@ is used.
    \baselineskip#2\p@
    \lineskip#3\ifx\UNDEFINED\ex@\p@\else\ex@\fi\relax
    \lineskiplimit\lineskip
    \if b#1\relax\box\z@\else\vss\fi
    \hbox to\dimen@{\hss\unhbox\tw@\hss}%
    \if t#1\relax\box\z@\else\vss\fi
  }% end \vbox
}% end \preced@

\newcommand{\tpx}[2]{\@actsymbol[#1][]p{#2}[]}
\newcommand{\tqx}[2]{\@actsymbol[#1][]q{#2}[]}
\newcommand{\ndx}[2]{\@actsymbol[#1][]d{#2}[]}
\newcommand{\lx}[1]{\@actsymbol[][]{\ell}{#1}[]}

\newcommand{\pureend}[2]{\ensuremath\actsymb{A}{#1:\firsttop{\angl{#2}}}}

\newcommand{\firsttop}[1]{\ensuremath\nthtop{1}{#1}}
\newcommand{\secondtop}[1]{\ensuremath\nthtop{2}{#1}}
\newcommand{\thirdtop}[1]{\ensuremath\nthtop{3}{#1}}
\newcommand{\firstbottom}[1]{\ensuremath\nthbottom{1}{#1}}
\newcommand{\secondbottom}[1]{\ensuremath\nthbottom{2}{#1}}
\newcommand{\thirdbottom}[1]{\ensuremath\nthbottom{3}{#1}}

\newcommand{\IA}{\ensuremath\twoletsymb{IA}}
\newcommand{\IAbar}{\ensuremath\twoletsymb{I\bar{A}}}
\newcommand{\IbarAbar}{\ensuremath\twoletsymb{\bar{I}\bar{A}}}
\newcommand{\DA}{\ensuremath\twoletsymb{DA}}
\newcommand{\DAbar}{\ensuremath\twoletsymb{D\bar{A}}}
\newcommand{\DbarAbar}{\ensuremath\twoletsymb{\bar{D}\bar{A}}}

\newcommand{\Ia}{\ensuremath\twoletsymb{Ia}}
\newcommand{\Ibarabar}{\ensuremath\twoletsymb{\bar{I}\bar{a}}}
\newcommand{\Iadd}{\ensuremath\twoletsymb{I\ddot{a}}}
\newcommand{\Is}{\ensuremath\twoletsymb{Is}}
\newcommand{\Isdd}{\ensuremath\twoletsymb{I\ddot{s}}}
\newcommand{\Da}{\ensuremath\twoletsymb{Da}}
\newcommand{\Dbarabar}{\ensuremath\twoletsymb{\bar{D}\bar{a}}}
\newcommand{\Dadd}{\ensuremath\twoletsymb{D\ddot{a}}}
\newcommand{\Ds}{\ensuremath\twoletsymb{Ds}}
\newcommand{\Dsdd}{\ensuremath\twoletsymb{D\ddot{s}}}

%\iffalse
%</package>
%\fi
%    \end{macrocode}
%
% Local Variables:
% mode: doctex
% TeX-master: t
% End:
